% Figure: imbalance force F = m e omega^2 versus rotation speed, on log-log
% style linear axes, for several eccentricities, with the omega^2 growth.
\begin{figure}[htbp]
\centering
\begin{tikzpicture}
\begin{axis}[
  width=12cm, height=7cm,
  xlabel={rotation speed $N$ (revolutions per minute)},
  ylabel={imbalance force $F_{\text{imb}}$ (\si{\newton})},
  xmin=0, xmax=6000, ymin=0, ymax=2000,
  axis lines=left,
  domain=0:6000, samples=120,
  legend pos=north west,
  legend cell align=left,
  grid=major, grid style={enggray!20},
  tick label style={font=\scriptsize},
  label style={font=\small},
  legend style={font=\scriptsize},
]
% F = m e omega^2, omega = N 2pi/60. m = 20 kg rotor.
% F = m e (2 pi N /60)^2 = m e (0.10472 N)^2 = m e 0.010966 N^2.
% e = 0.05 mm = 5e-5 m: coeff = 20*5e-5*0.010966 = 1.0966e-5
\addplot[engblue, very thick] {1.0966e-5 * x^2};
\addlegendentry{$e=0.05\,\si{\milli\meter}$ (balanced)}
% e = 0.15 mm: coeff = 20*1.5e-4*0.010966 = 3.290e-5
\addplot[engamber, very thick] {3.290e-5 * x^2};
\addlegendentry{$e=0.15\,\si{\milli\meter}$}
% e = 0.50 mm: coeff = 20*5e-4*0.010966 = 1.0966e-4
\addplot[engred, very thick] {1.0966e-4 * x^2};
\addlegendentry{$e=0.50\,\si{\milli\meter}$ (out of spec)}
\end{axis}
\end{tikzpicture}
\caption[Imbalance force versus rotation speed]{Rotating-imbalance
force $F_{\text{imb}}=m e\,\omega^{2}$ for a $20\,\si{\kilogram}$
rotor at three residual eccentricities, against rotation speed in
revolutions per minute. The quadratic growth in speed is the load
designer's central fact: tripling the speed raises the force ninefold
at fixed eccentricity, and the residual eccentricity allowed by an
ISO~1940/1 balance grade therefore tightens as the rated speed rises.
A rotor balanced to $e=0.05\,\si{\milli\meter}$ (blue) generates a
modest force even at $6{,}000$ revolutions per minute; the same rotor
left at $e=0.50\,\si{\milli\meter}$ (red) drives nearly
$2{,}000\,\si{\newton}$ of rotating bearing load, cyclic at the
rotation frequency, into the bearings, shaft, and foundation.}
\label{fig:vol03ch04:imbalance-force}
\end{figure}
